% Keep the heading, its lead sentence, the first subheading and the opening
% entries together: a References heading stranded at the foot of a page is a
% review defect in both editions, which paginate differently.
\Needspace{14\baselineskip}
\section*{References}

Only sources actually used above, with loci sufficient to check the claims
made from them.

\subsection*{Liturgical witnesses}

\begin{itemize}[leftmargin=1.2em, itemsep=0.15em]
\item \work{Missale Romanum}, editio typica (Vatican City: Typis Polyglottis
Vaticanis, 1962), pp.\ 392--393, marginal nos.\ 1552--1561. Facsimile:
\url{https://media.churchmusicassociation.org/pdf/missale62.pdf}, PDF pp.\
473--474. \textit{Controlling text for every Latin form printed above.}
\item \work{Missale Romanum}, editio iuxta typicam (New York: Benziger
Brothers, 1962), pp.\ 386--388. Page images in Internet Archive item
\texttt{MissaleRomanum1962RomanMissalColorLatin}, leaves
\texttt{n462}--\texttt{n464}:
\url{https://archive.org/details/MissaleRomanum1962RomanMissalColorLatin}.
\textit{Second 1962 edition, collated for the edition differences.}
\item \work{The Roman Missal translated into the English language for the use
of the laity}, first revised edition (Philadelphia: Eugene Cummiskey, 1861),
Collect, Secret and Postcommunion of this formulary, printed pp.\ 420, 422
and 423. Anonymous translation; public domain in the United States.
\textit{Source of the English of all three orations.} Registered as
\nolinkurl{edition.eugene-cummiskey.roman-missal-english-laity.philadelphia-1861},
passage \texttt{post-pentecosten-12}.
\item \work{Sacramentarium Leonianum}, ed.\ C.~L.~Feltoe (Cambridge:
University Press, 1896), pp.\ 60 and 74 (Masses of July).
\url{https://archive.org/details/sacramentariumle00cath}. \textit{Earliest
located witnesses to the Collect and Postcommunion.}
\item \work{The Gelasian Sacramentary}, ed.\ H.~A.~Wilson (Oxford: Clarendon,
1894), Book III, no.\ VIII, pp.\ 228--229, with the incipit index and the
appendix collation at p.\ 354.
\url{https://archive.org/details/gelasiansacrame00gelagoog}. \textit{All three
orations in one Mass-set, with the \latin{largiaris} reading.}
\item \work{The Gregorian Sacramentary under Charles the Great}, ed.\
H.~A.~Wilson, Henry Bradshaw Society 49 (London, 1915), no.\ LXXXI, p.\ 173;
the apparatus to no.\ LXXX at pp.\ 172--173 for the chant cues; and p.\ 199.
\url{https://archive.org/details/gregoriansacrame00cath}
\item \work{Missale Romanum}, editio typica tertia (Vatican City, 2008): the
Collect at \latin{Dominica XXXI per annum}, the Secret as \latin{super oblata}
of \latin{Feria quinta post Cineres}, the Postcommunion of \latin{Feria tertia
hebdomadae III Quadragesimae}. \textit{Cited for those assignments and, once,
for the official English of the super oblata; that translation is under
copyright and the single clause quoted does not substitute for the book.}
\item Karl Ott, \work{Offertoriale sive versus offertoriorum} (Tournai:
Desclée, 1935), no.\ 59, pp.\ 97--100.
\url{https://media.churchmusicassociation.org/books/offertoriale1935.pdf}.
\textit{The two historic offertory verses, no longer appointed.}
\item Dominic Johner, \work{The Chants of the Vatican Gradual} (Collegeville:
St.\ John's Abbey Press, 1940), pp.\ 296--297.
\url{https://archive.org/details/chantsofvaticang00john}
\item Ildefonso Schuster, \work{The Sacramentary (Liber Sacramentorum)},
English edition (London: Burns Oates \& Washbourne, 1924--30), vol.~III,
``Twelfth Sunday after Pentecost,'' pp.\ 128--132.
\url{https://archive.org/details/LiberSacramentorum}. \textit{The English
translation may still be in copyright in some jurisdictions; quoted only in
short compass.}
\end{itemize}

\subsection*{Scripture}

\begin{itemize}[leftmargin=1.2em, itemsep=0.15em]
\item \work{The Holy Bible, Douay--Rheims}, the Challoner revision
(1749--1752), in the Project Gutenberg e-text of that revision (eBook 1581).
Quoted for Ps.\ 33:2--3; 69:2--4; 87:2; 103:13--15; Ex.\ 32:11--14; Lk.\
10:23--37; 2 Cor.\ 3:4--9. \textit{Source of the English of every scriptural
proper.} Public domain by age. Registered as
\nolinkurl{edition.english-college-of-douay.douay-rheims-bible.challoner-gutenberg-1581},
with its book index, Vulgate/Hebrew/English psalm-numbering concordance,
verse-alias table, and collation against the 1899 American edition. Exactly
one quoted locus appears in that collation: Ex.\ 32:11, where the e-text reads
``is thy indignation enkindled'' and the American edition prints ``kindled.''
\item Clementine Vulgate, consulted for Ps.\ 33, 39, 69, 87, 103; Ex.\ 3, 32,
33; Lk.\ 10; 2 Cor.\ 3. \url{https://www.drbo.org/lvb/}
\end{itemize}

\subsection*{Patristic, medieval and later reception}

\begin{itemize}[leftmargin=1.2em, itemsep=0.15em]
\item Ambrose, \work{Expositio evangelii secundum Lucam} VII, §§69--84 (PL 15,
coll.\ 1717--1720), and the lemma scan of Book VII that establishes his
silence at Lk.\ 10:23--24.
\url{https://la.wikisource.org/wiki/Expositio_Evangelii_secundum_Lucam_(Ambrosius)/7}
\item Ambrosiaster, \work{Commentaria in Epistolam ad Corinthios Secundam}, ad
3:6--9 (PL 17; the source page displays no column numbers).
\url{https://la.wikisource.org/wiki/Commentaria_in_Epistolam_ad_Corinthios_Secundam_(Ambrosiaster)}
\item Augustine, \work{Enarrationes in Psalmos} 33 (two sermons; New Advent's
``Psalm 34'' carries only the second, so Sermo~I is quoted from the Latin);
69 §§1--4, 7 (New Advent ``Psalm 70''); 87 §§1--2 (``Psalm 88''); 103 Sermo~3
§§12--14 (``Psalm 104'' §§19--21). English in NPNF first series, vol.~8:
\url{https://www.newadvent.org/fathers/1801.htm}. Latin:
\url{https://www.augustinus.it/latino/esposizioni_salmi/}
\item Augustine, \work{Quaestiones Evangeliorum} II.19; \work{Sermo} 131.6;
\work{De doctrina christiana} I.30 and III.5.9; \work{De spiritu et littera}
4.6, 14.23--25 and 19.34; \work{Quaestiones in Heptateuchum} II, qq.\ 10 and
141--149; \work{De civitate Dei} XV.25. English in NPNF first series, vols.\ 2
and 5 (\url{https://www.newadvent.org/fathers/}); Latin at
\url{https://www.augustinus.it/latino/} and
\url{https://www.thelatinlibrary.com/augustine/}
\item Basil the Great, \work{Homilia in Psalmum 33} (PG 29, coll.\ 349--356),
read in the public Migne scan with the parallel Latin. No public-domain English
of this homily exists; the copyrighted Fathers of the Church translation was
not used.
\item Bede, \work{In Lucae Evangelium Expositio} III, on Lk.\ 10:23--37
(PL 92, coll.\ 467C--470D). Registered in the repository's source library as
passage
\nolinkurl{passage.bede.in-lucae-evangelium-expositio.latin-migne-corpus-corporum-web-2026-07-25.iii-lk-10.23-37}.
\item Robert Bellarmine, \work{Explanatio in Psalmos}, Pss.\ XXXIII, LXIX,
LXXXVII and CIII, tr.\ John O'Sullivan (Dublin: James Duffy, 1866), artifact
pp.\ 100--105, 220, 275--278 and 333--339.
\url{https://archive.org/details/commentaryonbook0000bell}
\item Bernard of Clairvaux, \work{Sermones super Cantica Canticorum} 7.5
(PL 183, col.\ 809B).
\item John Cassian, \work{Conferences} X.10. NPNF second series, vol.~11:
\url{https://www.newadvent.org/fathers/350810.htm}
\item Cassiodorus, \work{Expositio Psalmorum} 33, 69 (PL 70, coll.\ 493--494),
87 and 103, in the public Migne text. The copyrighted ACW translation was not
used.
\item John Chrysostom, \work{Homilies on Second Corinthians} 6 and 7;
\work{Homily 45 on Matthew}; \work{Homily 16 on Romans}; and \work{Homily 10
on Hebrews} as excerpted in the \work{Catena aurea}. NPNF first series,
vols.\ 10, 11 and 12: \url{https://www.newadvent.org/fathers/220206.htm},
\url{https://www.newadvent.org/fathers/220207.htm},
\url{https://www.newadvent.org/fathers/210216.htm}
\item Clement of Alexandria, \work{Quis dives salvetur} 28--29. ANF vol.~2:
\url{https://www.newadvent.org/fathers/0207.htm}
\item Cyril of Alexandria, \work{Commentary on Luke}, Sermons 67 and 68, tr.\
R.~Payne Smith (Oxford, 1859).
\url{https://www.tertullian.org/fathers/cyril_on_luke_09_sermons_63_69.htm}
\item Gregory the Great, \work{Moralia in Iob} IX.xvi.23 and XX.v.14;
\work{Regula pastoralis} II.6; \work{Homiliae in Evangelia} II.18.2, and the
enumeration of all forty homilies that establishes the bounded negative.
Library of the Fathers translation:
\url{https://www.lectionarycentral.com/GregoryMoralia/Book09.html},
\url{https://www.lectionarycentral.com/GregoryMoralia/Book20.html}
\item Irenaeus, \work{Adversus haereses} III.17.3. ANF vol.~1:
\url{https://www.newadvent.org/fathers/0103317.htm}
\item Jerome, \work{Epistle} 108.12 (Adummim, ``the Place of Blood''). NPNF
second series, vol.~6: \url{https://www.newadvent.org/fathers/3001108.htm}
\item Origen, \work{Homily 34 on Luke} in Jerome's Latin (PL 26); \work{De
principiis} I.1.2, with the bounded negative at Book IV; \work{Contra Celsum}
VII.20. Latin:
\url{https://la.wikisource.org/wiki/Translatio_XXXIX_Homiliarum}. English
(ANF vol.~4): \url{https://www.newadvent.org/fathers/04121.htm},
\url{https://www.newadvent.org/fathers/04167.htm}. The modern English of the
Lucan homilies (FC 94) is in copyright and is not quoted.
\item \work{The Rule of St.\ Benedict}, cc.\ 17--18 and 35. Latin:
\url{https://www.thelatinlibrary.com/benedict.html}. English, Boniface
Verheyen (1902): \url{https://www.ccel.org/ccel/benedict/rule.html}
\item Tertullian, \work{Adversus Marcionem} II.26. ANF vol.~3:
\url{https://www.tertullian.org/anf/anf03/anf03-29.htm}
\item Theodoret of Cyrus, \work{Interpretatio in Psalmos} 69 (PG 80, coll.\
1415--1418), 33 (coll.\ 1101--1109), 87 (coll.\ 1567--1573, the last figure
approximate) and 103 (coll.\ 1693--1707), read on the registered facsimile of
PG 80 and quoted from Migne's parallel Latin; \work{Interpretatio} in
2 Corinthians (PG 82, from col.\ 376), verified as publicly reachable but not
quoted. Robert Hill's modern English translations are in copyright and were
not used.
\item Thomas Aquinas, \work{Super II ad Corinthios}, cap.~3, lect.\ 1--2
(\url{https://www.corpusthomisticum.org/c2c.html}); \work{Summa theologiae}
I-II q.~106 aa.\ 1--2
(\url{https://www.corpusthomisticum.org/sth2106.html}; English, public domain,
\url{https://www.newadvent.org/summa/2106.htm}); \work{Catena aurea} on Lk.\
10, lect.\ 7--9, in the Oxford translation of the 1840s, whose marginal
sigla and preface carry the pseudo-Chrysostom finding. Larcher's English of
the Pauline commentary is in copyright and was used only for checking.
\item John Calvin, \work{Harmony of the Evangelists}, at Lk.\ 10:30, tr.\
William Pringle. \url{https://ccel.org/ccel/calvin/calcom32}
\end{itemize}

\subsection*{Magisterial and official texts}

\begin{itemize}[leftmargin=1.2em, itemsep=0.15em]
\item John Paul II, \work{Salvifici doloris} (11 February 1984), §§28--30.
\url{https://www.vatican.va/content/john-paul-ii/en/apost_letters/documents/hf_jp-ii_apl_11021984_salvifici-doloris.html}
\item Pontifical Biblical Commission, responses of 1906 on the Mosaic
authorship of the Pentateuch, as reported in \work{The Catholic
Encyclopedia}, ``Pentateuch'' (vol.~11, 1911).
\url{https://www.newadvent.org/cathen/11646c.htm}
\end{itemize}

\subsection*{Historical orientation and cultural afterlives}

\begin{itemize}[leftmargin=1.2em, itemsep=0.15em]
\item \work{The Catholic Encyclopedia} (New York, 1907--1912), all at
\url{https://www.newadvent.org/cathen/}: ``Breviary'' (vol.~2, 1907), for the
\work{Psalterium Romanum} still in liturgical use at Rome; ``Epistles to the
Corinthians'' (vol.~4, 1908); ``Golden Calf'' (vol.~6, 1909); ``Gospel of
Saint Luke'' (vol.~9, 1910); ``Offertory'' (vol.~11, 1911); ``Parables''
(vol.~11, 1911); ``Pentateuch'' (vol.~11, 1911); and ``Psalms'' (vol.~12,
1911).
\item \work{New American Bible, Revised Edition}, introductions to Luke and to
2 Corinthians and the note to Ps.\ 70, at
\url{https://bible.usccb.org/bible/}. Summarised; protected text is not
reproduced.
\item \work{Donoghue v Stevenson} [1932] A.C.\ 562 at 580 (House of Lords, 26
May 1932), Lord Atkin. Read in the transcription of the judgment published by
the Scottish Council of Law Reporting, Lord Atkin page 2:
\url{https://www.scottishlawreports.org.uk/resources/donoghue-v-stevenson/appeal-papers/index-of-judgement/lord-atkin-page-1/}
\item John M.\ Darley and C.\ Daniel Batson, ```From Jerusalem to Jericho': A
Study of Situational and Dispositional Variables in Helping Behavior,''
\work{Journal of Personality and Social Psychology} 27.1 (1973), pp.\
100--108. Under APA copyright; only short sentences are quoted, with citation.
\item \work{United States ex rel.\ Knauff v.\ Shaughnessy}, 338 U.S.\ 537
(1950), Frankfurter J., dissenting, at 548 and 550, in the official United
States Reports. Public domain.
\item Claudio Monteverdi, \work{Vespro della Beata Vergine} (Venice: Ricciardo
Amadino, 1610), opening versicle and response; and Jennifer More Glagov,
programme note, Music of the Baroque (2016), for the documented borrowing of
the \work{Orfeo} toccata.
\item Giovanni Pierluigi da Palestrina, \work{Precatus est Moyses}, 5vv, in
\work{Offertoria totius anni} (Rome: Francesco Coattino, 1593), under the
rubric \latin{Dominica XII post Pentecosten}; Orlande de Lassus,
\work{Precatus est Moyses}, LV 858, in \work{Sämtliche Werke} vol.~3
(Leipzig: Breitkopf \& Härtel, 1895), pp.\ 23--25. \textit{Cited for their
existence and publication facts only.}
\end{itemize}
